\section{引言}

格点计算的目标之一是仅从 QCD 拉格朗日量出发，
预言禁闭夸克与胶子的低能强子谱。这一谱学进路必须有测量具有所研究
量子数的场算符两点关联函数的手段。对 QCD 谱的完整理解
还必须包括介子与重子的激发态以及奇特态。
把散射实验中看到的共振态描述为 QCD 中的态，长期以来对格点计算是一大挑战，
因为在欧氏场论表述中通常无法直接获得与衰变宽度相联系的矩阵元。
规避这一困难的途径原则上众所周知：可以通过细致研究有限盒子中谱
对盒子体积的依赖性来推断衰变性质 \cite{DeWitt:1956be,Luscher:1991cf}。
对能级做出可靠确定的最佳手段是采用庞大的算符基，
然后用变分法寻找能量本征态的良好近似。一旦要探索阈值以上的态，
算符基就应包含那些类似多强子体系、且组分具有确定动量的算符。
获取与长程关联相关的夸克传播子的全部元素 \cite{Foley:2005ac} 使这些测量成为可能，
也使人们得以更细致地研究那些其精确测量长期困扰格点工作者的物理可观测量。
一个例子是同位旋标量介子领域：其质量测定需要在两点关联函数中
计算不连通图。除了这份雄心勃勃的要求清单之外，高精度数据
对这些计算也将至关重要 \cite{Morningstar:2008mc}。

在格点上的蒙特卡罗计算中，关联函数中物理相关的信号指数式衰减，
并迅速被统计涨落淹没。能迅速产生低能能量本征态的算符极其宝贵，
它们能指数级地改善所提取数据的质量。格点工作者手中
构建良好产生算符最有用的工具是涂抹。因此，理解禁闭低能自由度的最佳途径
必须来自一种能兼顾所有这些特征的涂抹方法：它在数值上必须易于实现，
必须能很好地投影到最低能量的态上，并且必须便于计算涉及广泛的
产生算符集合（包括激发多强子态的那些算符）的关联函数。

新一代致力于强子谱学的实验，包括 Jefferson Lab 的 GlueX、
GSI/FAIR 的 PANDA 以及 BES III，打算以前所未有的精度、
在此前未探索过的质量区间与量子数领域开展测量。Hadron Spectrum
合作组的主要目标是从现实的格点模拟中提取这些领域中
质量、衰变性质及相关矩阵元的预言。

谱学研究几乎与非微扰蒙特卡罗格点计算一样古老
\cite{Hamber:1981zn,Weingarten:1980hx}。随着最初研究的推进，
人们很快意识到：直接用格点拉格朗日量中的场构造的产生强子的算符
与一整座态的``高塔''有显著重叠，从而难以提取谱中最轻的部分。
解决办法是用``涂抹''过的场来构造算符：在形成产生算符之前，
先把某个线性算符作用于适当时间片上的夸克自由度。
构造涂抹场的目的在于大幅压低接近截断尺度的场分量，
因为这些模式对长程关联函数的贡献不大。
最初的技术考虑的是非规范协变的涂抹场。不久之后便引入了
规范协变方案 \cite{Gusken:1989ad,
Allton:1993wc}。特别地，文献~\cite{Allton:1993wc} 提出了 Jacobi
涂抹算法：先构造三维规范协变拉普拉斯算符（laplacian）的格点近似，
再把它迭代地作用于场。这自然起到了长波滤波器的作用，
过滤构成夸克场的各个模式。


本文提出了一种新的夸克场涂抹方法，并在大规模现实模拟中进行了检验。
第~\ref{sec:theory} 节给出新算法的系统表述，
包括其在强子两点关联函数以及插入任意流的三点函数测量中的应用细节。
第~\ref{sec:results} 节给出该方法有效性的初步检验结果；
第~\ref{sec:discussion} 节讨论该方法引发的实际问题，
并概述若干未来研究方向。
